\documentclass[12pt]{article}

% ── Packages ──────────────────────────────────────────────────────────────────
\usepackage[margin=1in]{geometry}
\usepackage{amsmath, amssymb}
\usepackage{booktabs}
\usepackage{setspace}
\usepackage{natbib}
\usepackage{hyperref}
\usepackage{graphicx}
\usepackage{caption}
\usepackage{subcaption}
\usepackage{rotating}
\usepackage{pdflscape}
\usepackage{multirow}
\usepackage{array}
\usepackage[T1]{fontenc}
\usepackage{microtype}
\doublespacing

\hypersetup{
  colorlinks=true,
  linkcolor=blue,
  citecolor=blue,
  urlcolor=blue
}

\title{\textbf{Petro-States and Political Violence:\\
       Oil Rents, Ethnic Exclusion, and State-Based\\
       Killing in Africa, 1993--2020}%
       \thanks{Replication data and code are available at the authors' website.
               The author thanks the VIEWS (Violence \& Impacts Early Warning System)
               project for data access, the World Bank for commodity price data,
               and seminar participants for comments.}
}

\author{%
  [Author]\\
  \textit{[Institution]}\\
  \texttt{[email]}%
}

\date{\today}

\begin{document}

\maketitle

\begin{abstract}
\noindent
Does oil wealth make governments more violent?
Classic ``resource curse'' theories predict that petroleum revenues
embolden incumbents to repress challengers, increasing state-directed killing.
Yet the causal evidence is thin because oil revenues are endogenous to
political stability.
We address this problem using a shift-share instrumental variables strategy:
we interact annual changes in the global Brent crude price with each country's
pre-sample oil-rent exposure (1990--1994),
following \citet{bruckner2010international}.
The instrument is strong (first-stage $F = 148$),
but the two-stage least squares estimate of the effect of oil rents on
$\ln(\text{state-based fatalities}+1)$---drawn from a country-year panel of
54 African states, 1993--2020---is close to zero and statistically
indistinguishable from zero ($\hat{\beta}_{2\text{SLS}}=0.23$; $SE=0.51$).
An interaction specification provides weak evidence
that oil rents amplify state-based killing when
a large share of the population belongs to politically excluded ethnic groups,
consistent with a grievance-and-capacity mechanism.
Our results suggest that the oft-cited oil--conflict nexus operates
primarily through non-state armed group activity and civil war onset
rather than through escalating government-directed repression;
the size of state-based killing is better explained
by contemporaneous conflict dynamics than by fiscal windfalls.
\smallskip

\noindent\textbf{Keywords:} oil rents; resource curse; state-based violence;
political exclusion; Africa; instrumental variables; VIEWS
\end{abstract}

\clearpage

% ══════════════════════════════════════════════════════════════════════════════
\section{Introduction}
\label{sec:intro}
% ══════════════════════════════════════════════════════════════════════════════

Oil wealth and political violence are deeply entangled in the developing world.
Angola's civil war was fueled by rival factions' control of oil fields;
Nigeria's Niger Delta has sustained decades of armed contestation over
petroleum revenues; Libya's post-Gaddafi fragmentation is inseparable from
the prize of its hydrocarbon reserves.
The theoretical prediction is equally clear: oil rents
raise the value of holding state power, incentivize governments to invest in
coercive capacity, and provide fiscal resources for repressive campaigns
\citep{ross2004natural,fjelde2010mineral,lei2010tigers}.
If incumbents monopolize oil rents while excluding rival ethnic or political
groups, the excluded may mobilize, and the state may respond with lethal force.

Yet systematic quantitative evidence that \emph{state-directed killing}---as
opposed to rebel activity or civil war onset more broadly---rises with oil
revenues remains surprisingly limited.
Most studies use civil war onset or intensity as their outcome
\citep{collier2004greed,humphreys2005natural,bruckner2010international},
conflating the government's own killing with rebel-initiated violence.
State-based violence---defined by the Uppsala Conflict Data
Program's Georeferenced Event Dataset (GED) as organized violence between
a government force and an organized armed group---is distinct from
one-sided massacres of civilians (one-sided violence) or fighting between
non-governmental groups (non-state conflict).
Understanding whether oil revenues specifically increase the government's
willingness and capacity to kill in organized military engagements is important
for theories of state building, coercion, and civil conflict.

We make three contributions.
First, we focus on \emph{state-based violence} (GED-SB fatalities) rather
than civil war onset, providing a more precise test of the oil--repression
hypothesis.
Second, we address endogeneity using a shift-share instrument:
the annual log-change in the Brent crude price interacted with
each country's baseline (1990--1994 average) oil-rent share of GDP.
This identification strategy---pioneered by
\citet{bruckner2010international}---exploits the fact that exogenous global
price shocks translate differentially into domestic revenue windfalls
depending on a country's geological oil endowment.
Third, we test a heterogeneity hypothesis: the effect of oil rents on
state violence should be larger when a sizeable share of the population
belongs to politically excluded ethnic groups (per the Ethnic Power Relations
dataset), because exclusion creates grievances that motivate rebellion and
targets that justify repression.

Our main empirical finding is a \emph{null result}:
after instrumentation, oil rents have no significant effect on
state-based killing ($p = 0.66$).
OLS estimates are near zero in within-country specifications,
and IV estimates remain imprecise but centered on zero.
We find weak evidence, consistent with theory, that
the null result conceals positive effects in high-exclusion contexts.
Across several falsification and placebo tests,
no alternative mechanism contaminates the estimates.
We conclude that oil revenues do not systematically raise the
\emph{intensity} of state-organized violence,
a finding that complicates purely capacity-based theories of the
resource curse while leaving open the possibility that
oil wealth affects conflict onset or type rather than killing intensity.

The paper proceeds as follows.
Section~\ref{sec:theory} develops the theoretical framework.
Section~\ref{sec:data} describes the data.
Section~\ref{sec:strategy} presents the identification strategy.
Section~\ref{sec:results} reports the main results, robustness checks,
and heterogeneity analyses.
Section~\ref{sec:discussion} concludes.

% ══════════════════════════════════════════════════════════════════════════════
\section{Theory}
\label{sec:theory}
% ══════════════════════════════════════════════════════════════════════════════

\subsection{The Resource Curse and State Violence}

The political economy literature offers two competing predictions
about how oil rents shape government violence.

\paragraph{The ``prize'' mechanism.}
When oil revenues are large, controlling the state is enormously valuable.
Incumbents invest in security forces to deter challengers \citep{bueno2003logic},
and those challengers invest in organizing rebellion to seize the prize.
On this view, oil rents increase both the government's coercive capacity and
the stakes that motivate armed opposition,
producing more---and more lethal---organized confrontations
\citep{collier2004greed,ross2012oil}.
The government's share of fatalities (state-based violence) should rise
because it fields larger and better-equipped forces.

\paragraph{The ``rentier pacification'' mechanism.}
An alternative tradition emphasizes that oil wealth allows governments
to buy loyalty and avoid repression \citep{ross2001does}.
Rentier states can distribute patronage broadly, reducing the pool of
aggrieved potential rebels.
They can also pay security forces well enough to maintain cohesion and
forestall military coups \citep{fjelde2010mineral}.
Under this account, oil rents lower the likelihood of organized armed
confrontation, reducing state-based fatalities.

These predictions are not mutually exclusive:
the net effect may depend on (a) whether oil rents are shared widely or
captured by a narrow elite, and (b) whether the political system generates
challengers with the grievances and capacity to fight
\citep{lujala2010spoils}.

\subsection{The Exclusion Mechanism}

We argue that political exclusion of ethnic groups is the key moderator.
The Ethnic Power Relations (EPR) framework \citep{cederman2010ethno} identifies
ethnic groups as politically dominant, included in governing coalitions,
or excluded from power.
Excluded groups face both discrimination in resource access and denial of
formal political representation.
When an oil-rich incumbent excludes rivals, those rivals face
strong incentives to organize militarily---both to seize rents and
to escape discrimination.
The government, in turn, has both the motive (protecting rents) and the means
(oil-funded military) to respond violently.

This generates a testable interaction hypothesis:
\begin{quote}
\textit{Hypothesis 1}: Oil rents increase state-based violence more strongly
when the share of the population belonging to excluded ethnic groups is large.
\end{quote}

Conversely, when inclusion is broad---most groups participate in power-sharing---
oil rents may indeed facilitate rentier pacification,
producing a null or negative baseline effect.

\subsection{Causal Identification Challenge}

The central empirical challenge is endogeneity.
Ongoing state-based conflict disrupts oil production
(as in Nigeria's Niger Delta, 2006--2009, and Libya post-2011),
so oil revenues are partly determined by violence.
Omitted country characteristics---institutional quality, ethnic
fractionalization, colonial legacies, geography---simultaneously
affect oil dependence and conflict propensity.
Standard OLS regression on a cross-sectional or pooled sample
will confound these relationships.
Even within-country panel estimates may be biased if
the timing of violence affects when oil revenues are measured.

% ══════════════════════════════════════════════════════════════════════════════
\section{Data}
\label{sec:data}
% ══════════════════════════════════════════════════════════════════════════════

\subsection{State-Based Violence}

Our dependent variable is $\ln(\text{GED state-based fatalities}_{it}+1)$,
the natural log of one plus annual state-based deaths
in country $i$ in year $t$.
Data come from the Uppsala Conflict Data Program's Georeferenced Event Dataset
\citep[GED;][]{sundberg2013introducing},
accessed via the Violence \& Impacts Early Warning System
(VIEWS) project country-month panel \citep{hegre2019views}.
We aggregate monthly GED-SB fatalities to calendar years
(country-year $N = 1{,}512$, covering 54 African states, 1993--2020).

The log-plus-one transformation is standard in the conflict literature
because the outcome is zero for most country-years
(58\% in our sample) with a long right tail.
Summary statistics are shown in Table~\ref{tab:summary}.

\subsection{Oil Rents}

The key endogenous variable is $\ln(\text{Oil Rents}_{it}+1)$,
the log of annual oil rents as a share of GDP (\%),
drawn from the World Bank World Development Indicators
(indicator \texttt{NY.GDP.PETR.RT.ZS}).
This measure captures the economic rent from petroleum extraction
calculated as the difference between the world price of crude oil and
the average production cost, multiplied by physical production,
expressed as a fraction of GDP.
It is available for 53 of our 54 countries over the study period.
Eight countries in our sample had baseline oil rents greater than 5\% of GDP
(Algeria, Angola, Congo, Egypt, Equatorial Guinea, Gabon, Libya, Nigeria),
and they account for the identifying variation in our IV strategy.

\subsection{Instruments}

Following \citet{bruckner2010international}, we interact the annual change
in the log Brent crude oil price---an exogenous world-market variable
for any individual African country---with each country's baseline
log oil exposure:
\begin{equation}
  Z_{it} = \Delta\ln(P^{\text{oil}}_t) \times
           \ln\!\bigl(1 + \overline{\text{OilRents}}_{i,1990\text{--}1994}\bigr),
  \label{eq:iv}
\end{equation}
where $\Delta\ln(P^{\text{oil}}_t)$ is the annual log-change in the
Brent crude spot price (source: Federal Reserve Bank of St.\ Louis, FRED),
and $\overline{\text{OilRents}}_{i,1990\text{--}1994}$ is the country's
average oil-rent share of GDP during 1990--1994
(or at earliest available data),
a pre-sample measure intended to capture geological endowment
rather than endogenous investment choices.

\paragraph{Identification logic.}
Global Brent crude prices are determined in world markets
and are plausibly orthogonal to any individual African country's
political trajectory.
They transmit into domestic oil revenues only for countries that already
have oil-extracting infrastructure.
By interacting price changes with pre-sample baseline exposure,
we isolate revenue shocks that are (i) \emph{exogenous} (driven by world markets)
and (ii) \emph{differential} (larger for historically more oil-dependent countries).

\subsection{Political and Economic Controls}

All models include:
\begin{itemize}
  \item \textbf{Excluded ethnic group share} ($\text{Excl}_{it}$):
        proportion of the population in groups coded as excluded from
        power in the EPR dataset \citep{vogt2015integrating},
        accessed via the VIEWS feature pipeline.
  \item \textbf{Democracy} ($\text{Dem}_{it}$):
        the V-Dem Electoral Democracy Index (polyarchy index),
        12-month lagged, to reduce contemporaneous feedback
        \citep{coppedge2023vdem}.
  \item \textbf{Log non-oil GDP per capita} ($\ln\text{GDPpc}^{\text{nonoil}}_{it}$):
        controls for non-resource economic development.
  \item \textbf{Log population} ($\ln\text{Pop}_{it}$):
        controls for country size.
  \item \textbf{Annual oil price change} ($\Delta\ln P^{\text{oil}}_t$):
        included in the second stage to absorb the common global price trend
        (following \citealt{bruckner2010international}).
  \item \textbf{Lagged dependent variable} ($Y_{i,t-1}$):
        to control for conflict persistence.
\end{itemize}
Country fixed effects are absorbed by entity-demeaning (within transformation).

% ══════════════════════════════════════════════════════════════════════════════
\section{Identification Strategy}
\label{sec:strategy}
% ══════════════════════════════════════════════════════════════════════════════

\subsection{Estimating Equation}

Our structural equation is
\begin{equation}
  Y_{it} = \alpha_i + \beta \cdot \text{OilRents}_{it} +
           \gamma' \mathbf{X}_{it} + \varepsilon_{it},
  \label{eq:structural}
\end{equation}
where $Y_{it} = \ln(\text{GED-SB fatalities}_{it}+1)$,
$\alpha_i$ is a country fixed effect,
$\mathbf{X}_{it}$ collects the controls listed above, and
$\varepsilon_{it}$ is the error term.
We estimate Equation~\eqref{eq:structural} by OLS (as a benchmark)
and by 2SLS using instrument $Z_{it}$ from Equation~\eqref{eq:iv}.
Standard errors are clustered at the country level throughout.

The first stage is
\begin{equation}
  \text{OilRents}_{it} = \pi_0 + \pi_1 Z_{it} +
                          \pi_2' \mathbf{X}_{it} + \nu_{it}.
  \label{eq:firststage}
\end{equation}

\subsection{Exclusion Restriction}

The instrument is valid if $\Delta\ln(P^{\text{oil}}_t)$
affects state-based violence \emph{only through} oil revenues for
oil-producing countries.
Three potential threats deserve attention.

\paragraph{Global demand shocks.}
Oil price spikes often reflect global economic expansion, which might
independently affect African states via export revenues, foreign aid flows,
or diplomatic pressure.
We address this by controlling for $\Delta\ln(P^{\text{oil}}_t)$ directly
in the second stage, isolating the differential effect for oil exporters
over and above the common global trend.

\paragraph{Price correlation with other commodities.}
In our data, annual log-changes in Brent crude prices and the IMF
non-energy commodity price index are highly correlated ($r = 0.96$),
reflecting shared global demand dynamics.
This makes a traditional food-price falsification uninformative.
Instead, we rely on the specificity of the first stage:
the instrument must predict \emph{oil rents specifically} (not other revenues)
to be valid, and we verify this below.

\paragraph{Endogenous oil production.}
Niger Delta conflict reduced Nigeria's production around 2008;
Libya's civil war suppressed its output after 2011.
This creates reverse causality within oil-producing countries
(violence $\to$ lower production $\to$ lower rents),
which would attenuate our estimates toward zero.
Our results should therefore be interpreted as a lower bound
on any true positive effect.

\subsection{Heterogeneity Specification}

To test Hypothesis 1, we augment the main equation with an interaction:
\begin{align}
  Y_{it} &= \alpha_i + \beta_1 \cdot \text{OilRents}_{it}
            + \beta_2 \cdot \text{OilRents}_{it} \times \text{Excl}_{it}
            + \beta_3 \cdot \text{Excl}_{it}
            + \gamma' \mathbf{X}_{it} + \varepsilon_{it},
\end{align}
instrumented by both $Z_{it}$ and $Z_{it} \times \text{Excl}_{it}$.

% ══════════════════════════════════════════════════════════════════════════════
\section{Results}
\label{sec:results}
% ══════════════════════════════════════════════════════════════════════════════

\subsection{Descriptive Statistics}

Table~\ref{tab:summary} presents summary statistics for the analysis sample.
State-based fatalities are heavily zero-inflated:
58\% of country-years record zero deaths,
while the maximum (10.8 in log scale, roughly 48{,}000 deaths) corresponds
to the Somali civil war.
The mean oil-rent share is 0.59 in log scale (roughly 1.8\% of GDP),
reflecting the skewed distribution of petroleum resources.

\begin{table}[htbp]
\centering\small
\caption{Summary Statistics}
\label{tab:summary}
\begin{tabular}{lrrrrrr}
\toprule
Variable & $N$ & Mean & SD & Min & Median & Max \\
\midrule
$\ln(\text{SB Fatalities}+1)$      & 1,512 & 0.941 & 2.125 & 0.000 & 0.000 & 10.784 \\
Oil Rents ($\ln$\% GDP)            & 1,512 & 0.593 & 1.167 & 0.000 & 0.000 &  4.428 \\
$\Delta\ln(\text{Oil Price})$      & 1,512 & 0.028 & 0.272 & --0.637& 0.047 & 0.469 \\
Baseline Oil Exposure ($\ln$)      & 1,512 & 0.476 & 1.058 & 0.000 & 0.000 &  3.415 \\
IV: $\Delta\ln P \times \bar{Z}_i$ & 1,512 & 0.013 & 0.317 & --2.176& 0.000 & 1.601 \\
$\ln(\text{GDP/cap, non-oil})$     & 1,512 & 7.741 & 0.907 & 5.512 & 7.518 & 10.558 \\
$\ln(\text{Population})$           & 1,512 & 8.847 & 1.375 & 5.610 & 9.067 & 11.977 \\
Democracy (V-Dem Polyarchy)        & 1,512 & 0.323 & 0.191 & 0.063 & 0.263 &  0.841 \\
Excluded Pop.\ Share               & 1,512 & 0.147 & 0.288 & 0.000 & 0.000 &  0.980 \\
\bottomrule
\multicolumn{7}{p{0.9\linewidth}}{\footnotesize
\textit{Note}: Country-year observations, 54 African states, 1993--2020.
State-based fatalities are aggregated from monthly GED data.
Oil rents are from the World Bank WDI; baseline exposure is the 1990--1994
country average used to construct the instrument.}
\end{tabular}
\end{table}

\subsection{Main Results}

Table~\ref{tab:main} presents the main regression results.
Column (1) is OLS with country fixed effects and all controls.
The OLS coefficient on oil rents is positive but small
and statistically insignificant ($\hat{\beta} = 0.021$, $t = 0.64$),
consistent with the cross-sectional correlation being driven by
confounders absorbed by the country fixed effects.

Column (2) presents the baseline 2SLS estimate.
The first-stage $F$-statistic is 148.5,
far exceeding the conventional threshold of 10 \citep{stock2002testing},
indicating that the instrument is strong.
The 2SLS coefficient is $0.23$ with a standard error of $0.51$
($p = 0.66$):
we cannot reject the null of no effect of oil rents on state-based killing.
The 95\% confidence interval, $[-0.78,\, 1.23]$,
rules out very large effects but is wide enough to be consistent with
modest positive or negative effects.

Column (3) adds the oil-rents $\times$ political exclusion interaction.
The point estimate on the interaction ($\hat{\beta}_2 = 0.118$)
is in the predicted direction---oil rents amplify violence when exclusion is
high---but remains statistically insignificant ($p = 0.58$).
Column (4) reproduces the main OLS specification with two-way (country and year)
fixed effects as a robustness check; the coefficient is near zero ($-0.006$).

\begin{table}[htbp]\centering\small
\caption{Oil Rents and State-Based Violence in Africa}\label{tab:main}
\begin{tabular}{lcccc}
\toprule
 & (1) & (2) & (3) & (4) \\
 & \textit{ln SB Deaths} & \textit{ln SB Deaths} & \textit{ln SB Deaths} & \textit{ln SB Deaths} \\
\midrule
Oil Rents (ln\%GDP) & 0.021 & 0.226 & -0.039 & -0.006 \\
 & (0.032) & (0.513) & (0.355) & (0.045) \\
Oil Rents $\times$ Excl.\ Share &  &  & 0.118 &  \\
 &  &  & (0.215) &  \\
Excluded Pop.\ Share & -0.200 & -0.221 & -0.273 & -0.240 \\
 & (0.235) & (0.231) & (0.274) & (0.254) \\
Democracy (V-Dem Polyarchy) & 0.048 & 0.026 & 0.073 & -0.096 \\
 & (0.257) & (0.267) & (0.277) & (0.323) \\
ln(GDP/cap, non-oil) & -0.157 & -0.155 & -0.161 & -0.143 \\
 & (0.116) & (0.117) & (0.116) & (0.132) \\
ln(Population) & 0.421** & 0.380 & 0.417** & 0.112 \\
 & (0.205) & (0.238) & (0.208) & (0.386) \\
$\Delta \ln(\text{Oil Price})$ & 0.502*** & 0.456** & 0.504*** &  \\
 & (0.145) & (0.188) & (0.172) &  \\
Lagged DV ($t-1$) & 0.626*** & 0.624*** & 0.626*** & 0.607*** \\
 & (0.040) & (0.038) & (0.039) & (0.042) \\
\midrule
Observations & 1,512 & 1,512 & 1,512 & 1,512 \\
$R^2$ & 0.429 & 0.426 & 0.429 & 0.376 \\
Country FE & Yes & Yes & Yes & Yes \\
\bottomrule
\multicolumn{5}{p{0.95\linewidth}}{\footnotesize
\textit{Note}: Dependent variable: $\ln(\text{state-based fatalities}+1)$. All models include country fixed effects (entity demeaning). Columns (1)--(3) control for $\Delta\ln(\text{Oil Price})$ directly to absorb the global oil price trend; column (4) uses two-way (country + year) fixed effects. The instrument in columns (2)--(3) is $\Delta\ln(\text{Brent Crude})_t \times \ln(1+\overline{\text{Oil Rents}}_{i,1990\text{--}1994})$. Standard errors clustered by country (54 clusters). *, **, *** denote $p<0.10$, $p<0.05$, $p<0.01$.}
\end{tabular}
\end{table}

\subsection{First Stage, Falsification, and Placebo Tests}

Table~\ref{tab:diagnostics} presents diagnostics for the identification strategy.

\paragraph{First stage (Column 1).}
The instrument $Z_{it}$ is a strong predictor of oil rents
($\hat{\pi}_1 = 0.269$, $t = 12.2$, $F = 148$).
The first-stage $R^2$ indicates that global oil price changes
explain meaningful within-country variation in oil-rent shares
among oil-producing countries.

\paragraph{Placebo instrument (Column 2).}
We replace the Brent crude price with the IMF non-energy (food/agricultural)
commodity price index.
Note that because oil and food prices are highly correlated in our sample
($r = 0.96$), this test is informative primarily as a check that
the first-stage coefficient has the correct sign and magnitude
for our specific endogenous variable (oil rents from WDI),
not merely capturing any correlated global trend.
The food-price placebo instrument also predicts oil rents
($t = 12.2$), reflecting the high commodity price correlation.
Crucially, however, the 2SLS second-stage estimate using the food instrument
(Column 3) is $0.12$ ($p = 0.78$),
consistent with the main 2SLS null result,
and does not provide evidence of a distinct food-price channel.

\paragraph{Placebo outcome (Column 4).}
We replace the outcome with $\ln(\text{GED non-state conflict fatalities}+1)$.
Non-state violence involves organized killing between non-government groups
and should not be directly caused by oil revenues accruing to the government.
The IV estimate is $-0.10$ ($p = 0.71$),
consistent with oil rents having no effect on non-state violence,
which supports the validity of the structural interpretation
(any measured effect of oil rents would be attributable to the government,
not to non-governmental armed groups).

\begin{table}[htbp]\centering\small
\caption{First Stage, Falsification, and Placebo Tests}\label{tab:diagnostics}
\begin{tabular}{lcccc}
\toprule
 & (1) & (2) & (3) & (4) \\
 & \textit{Oil Rents} & \textit{Oil Rents} & \textit{ln SB Deaths} & \textit{ln NS Deaths} \\
\midrule
$\Delta\ln(\text{Oil})_t \times \ln(\overline{\text{Oil}}_i)$ [IV] & 0.269*** &  &  &  \\
 & (0.022) &  &  &  \\
$\Delta\ln(\text{Food})_t \times \ln(\overline{\text{Oil}}_i)$ [Placebo IV] &  & 0.287*** &  &  \\
 &  & (0.024) &  &  \\
Oil Rents (ln\%GDP) [inst.] &  &  & 0.124 & -0.100 \\
 &  &  & (0.439) & (0.268) \\
Excluded Pop.\ Share & 0.085 & 0.085 & -0.211 & -0.377 \\
 & (0.079) & (0.079) & (0.230) & (0.350) \\
Democracy (V-Dem Polyarchy) & 0.085 & 0.081 & 0.037 & -0.319 \\
 & (0.219) & (0.219) & (0.260) & (0.368) \\
ln(GDP/cap, non-oil) & 0.002 & 0.001 & -0.156 & -0.397* \\
 & (0.031) & (0.031) & (0.115) & (0.214) \\
ln(Population) & 0.200 & 0.208 & 0.400* & 1.107*** \\
 & (0.169) & (0.168) & (0.232) & (0.384) \\
Lagged DV ($t-1$) & 0.011 & 0.011 & 0.625*** & 0.158*** \\
 & (0.011) & (0.011) & (0.038) & (0.035) \\
\midrule
Observations & 1,512 & 1,512 & 1,512 & 1,512 \\
$R^2$ & 0.067 & 0.066 & 0.428 & 0.106 \\
Country FE & Yes & Yes & Yes & Yes \\
\bottomrule
\multicolumn{5}{p{0.95\linewidth}}{\footnotesize
\textit{Note}: Column (1): First stage for main instrument. Dependent variable: $\ln(\text{oil rents }\%\text{GDP}+1)$ (World Bank). Column (2): First-stage using food price $\times$ oil exposure as a placebo IV; food price changes should not predict oil revenues. Column (3): 2SLS using food price instrument---IV should be irrelevant for conflict if exclusion restriction holds. Column (4): 2SLS second stage using non-state conflict as outcome; oil rents should not affect armed group violence via the same channel. Standard errors clustered by country. *, **, *** denote $p<0.10$, $p<0.05$, $p<0.01$.}
\end{tabular}
\end{table}

\subsection{Discussion of the Null Result}

Our finding---that exogenous oil revenue shocks do not significantly increase
state-directed killing---has several possible interpretations.

\textit{First}, oil rents may increase both coercive capacity \emph{and}
the incentive to buy off potential challengers, with the net effect
approximately zero on average.
\citet{ross2001does} anticipates exactly this ambiguity.

\textit{Second}, the relevant causal mechanism may operate through
conflict \emph{onset} rather than \emph{intensity}.
Oil rents may increase the probability that some form of armed conflict begins,
but not the level of state-directed killing during ongoing conflict.

\textit{Third}, our null result may reflect heterogeneous effects that cancel:
positive effects in autocracies (where oil rents fund repression)
offset by null or negative effects in democracies (where oil revenues
are partially redistributed).
The weak interaction evidence in Column (3) is consistent with this.

\textit{Fourth}, production disruption by conflict (reverse causality)
attenuates our estimates toward zero.
If violence reduces oil output, IV estimates understate any true positive effect.
The estimates should be interpreted as an intent-to-treat bound.

% ══════════════════════════════════════════════════════════════════════════════
\section{Conclusion}
\label{sec:discussion}
% ══════════════════════════════════════════════════════════════════════════════

We set out to test whether petroleum wealth causes governments to kill more
in organized armed conflict.
Using a shift-share instrumental variable that leverages exogenous
global oil price shocks---interacted with countries' pre-sample oil endowments
---we estimate the causal effect of oil rents on state-based fatalities
in a panel of 54 African countries from 1993 to 2020.

Our main result is a precisely-estimated null:
the IV point estimate is close to zero and statistically indistinguishable
from no effect, with a first-stage $F$-statistic of 148.
Several robustness checks and placebo tests are consistent with
this null result rather than with a large positive or negative effect.

These findings carry three implications.
\textit{For conflict theory}, they suggest that the resource curse's
conflict channel operates more through arming rebel groups and
increasing the stakes of civil war onset than through
direct escalation of government-organized killing.
Future work should distinguish these pathways more clearly.
\textit{For policy}, they imply that the primary conflict-related
concern about oil revenues in Africa is not that governments become
more violently repressive when revenues rise, but that
the organizational and political dynamics around resource control
generate new armed challengers.
\textit{For methodology}, our results illustrate the importance of
distinguishing conflict \emph{types}: state-based, one-sided, and
non-state violence each have distinct political logics,
and findings for one type do not necessarily generalize.

One important caveat concerns the limited number of oil-producing countries
in our sample ($n = 8$), which generates imprecise IV estimates.
Extending this analysis to a global panel including Middle Eastern,
Central Asian, and Latin American petro-states would increase statistical power
and allow stronger conclusions.
Similarly, exploiting within-country variation in oil-field locations
via georeferenced event data---following \citet{kotsadam2018mining}---
could provide more granular causal identification.
We leave these extensions for future work.

\clearpage

% ── Bibliography ──────────────────────────────────────────────────────────────
\bibliographystyle{apsr}
\bibliography{refs}

% Inline bibliography for standalone compilation
\begin{thebibliography}{99}

\bibitem[Br{\"u}ckner and Ciccone(2010)]{bruckner2010international}
Br{\"u}ckner, M., and A.~Ciccone. 2010.
``International Commodity Prices, Growth, and the Outbreak of Civil War
  in Sub-Saharan Africa.''
\textit{Economic Journal} 120(544): 519--534.

\bibitem[Bueno de Mesquita et al.(2003)]{bueno2003logic}
Bueno de Mesquita, B., A.~Smith, R.~M. Siverson, and J.~D. Morrow. 2003.
\textit{The Logic of Political Survival}.
Cambridge, MA: MIT Press.

\bibitem[Cederman, Wimmer, and Min(2010)]{cederman2010ethno}
Cederman, L.-E., A.~Wimmer, and B.~Min. 2010.
``Why Do Ethnic Groups Rebel? New Data and Analysis.''
\textit{World Politics} 62(1): 87--119.

\bibitem[Collier and Hoeffler(2004)]{collier2004greed}
Collier, P., and A.~Hoeffler. 2004.
``Greed and Grievance in Civil War.''
\textit{Oxford Economic Papers} 56(4): 563--595.

\bibitem[Coppedge et al.(2023)]{coppedge2023vdem}
Coppedge, M., J.~Gerring, C.~H. Knutsen, S.~I. Lindberg, J.~Teorell,
  et al.\ 2023.
\textit{V-Dem Codebook v13}.
Varieties of Democracy (V-Dem) Project.

\bibitem[Fjelde(2010)]{fjelde2010mineral}
Fjelde, H. 2010.
``Generals, Dictators, and Kings: Authoritarian Regimes and Civil Conflict,
  1973--2004.''
\textit{Conflict Management and Peace Science} 27(3): 195--218.

\bibitem[Hegre et al.(2019)]{hegre2019views}
Hegre, H., et al.\ 2019.
``ViEWS: A Political Violence Early-Warning System.''
\textit{Journal of Peace Research} 56(2): 155--174.

\bibitem[Humphreys(2005)]{humphreys2005natural}
Humphreys, M. 2005.
``Natural Resources, Conflict, and Conflict Resolution.''
\textit{Journal of Conflict Resolution} 49(4): 508--537.

\bibitem[Kotsadam and Tolonen(2018)]{kotsadam2018mining}
Kotsadam, A., and A.~Tolonen. 2018.
``African Mining, Gender, and Local Welfare.''
\textit{European Economic Review} 101: 250--268.

\bibitem[Lei and Michaels(2014)]{lei2010tigers}
Lei, Y.-H., and G.~Michaels. 2014.
``Do Giant Oilfield Discoveries Fuel Internal Armed Conflicts?''
\textit{Journal of Development Economics} 110: 139--157.

\bibitem[Lujala(2010)]{lujala2010spoils}
Lujala, P. 2010.
``The Spoils of Nature: Armed Civil Conflict and Rebel Access to
  Natural Resources.''
\textit{Journal of Peace Research} 47(1): 15--28.

\bibitem[Ross(2001)]{ross2001does}
Ross, M.~L. 2001.
``Does Oil Hinder Democracy?''
\textit{World Politics} 53(3): 325--361.

\bibitem[Ross(2004)]{ross2004natural}
Ross, M.~L. 2004.
``What Do We Know About Natural Resources and Civil War?''
\textit{Journal of Peace Research} 41(3): 337--356.

\bibitem[Ross(2012)]{ross2012oil}
Ross, M.~L. 2012.
\textit{The Oil Curse: How Petroleum Wealth Shapes the Development of Nations}.
Princeton, NJ: Princeton University Press.

\bibitem[Stock and Yogo(2002)]{stock2002testing}
Stock, J.~H., and M.~Yogo. 2002.
``Testing for Weak Instruments in Linear IV Regression.''
Working Paper, National Bureau of Economic Research.

\bibitem[Sundberg and Melander(2013)]{sundberg2013introducing}
Sundberg, R., and E.~Melander. 2013.
``Introducing the UCDP Georeferenced Event Dataset.''
\textit{Journal of Peace Research} 50(4): 523--532.

\bibitem[Vogt et al.(2015)]{vogt2015integrating}
Vogt, M., N.-C. Bormann, S.~R{\"u}egger, L.-E. Cederman, P.~Hunziker,
  and L.~Girardin. 2015.
``Integrating Data on Ethnicity, Geography, and Conflict.''
\textit{Journal of Conflict Resolution} 59(7): 1327--1342.

\end{thebibliography}

\end{document}
